\section{Agent: retroalimentación de ciclo corto y alta frecuencia, y el LM como núcleo}

Después de hablar de Xingmian (星眠) y Dokie, la atención pasa naturalmente al criterio sobre Agent que hay detrás de Dokie. Esta sección divide el know-how de Zhang Yueguang (张月光) sobre Agent en dos partes: la ruta de producto y el núcleo técnico.

La segunda mitad de la entrevista entra en el tema de Agent. Zhang Yueguang ofrece dos know-how. Primero, no hay que ver al Agent solo como una herramienta para tareas largas y fuera de línea: cuanto más larga es la tarea, más inestable se vuelve; quizá convenga construir ciclos más cortos, más frecuentes y más cercanos a la retroalimentación del usuario. Segundo, al menos hoy, el LM sigue siendo el núcleo intelectual del Agent, porque es capaz de entender la intención del usuario, descomponer tareas, invocar herramientas y llamar a otros modelos.

\begin{figure}[H]
\centering
\includegraphics[width=0.96\textwidth]{agent-short-loop.png}
\caption{Ruta de retroalimentación de ciclo corto y alta frecuencia para productos de Agent.}
\end{figure}

\begin{knowledgebox}{Lectura de la figura: el ciclo corto y de alta frecuencia no significa poca capacidad, sino un producto controlable}
En la figura, la intención del usuario entra al núcleo LM; el LM invoca herramientas y modelos para generar productos, y el usuario vuelve a dar retroalimentación con alta frecuencia. En comparación con meter al Agent en una tarea larga y esperar el resultado final, el ciclo corto y de alta frecuencia facilita construir confianza, reduce la ansiedad por la latencia e integra el juicio del usuario en el ciclo cerrado.
\end{knowledgebox}

\subsection{Por qué el LM sigue siendo el núcleo}

Zhang Yueguang no está muy de acuerdo con la afirmación de que «el lenguaje no puede llegar al punto final de la inteligencia»; incluso considera que el lenguaje es la esencia de la inteligencia. Esta postura forma un contraste interesante con la visión de los modelos del mundo de Xie Saining (谢赛宁) en EP133. Zhang Yueguang habla desde la perspectiva del producto y del Agent: la intención del usuario se expresa principalmente con lenguaje, la invocación de herramientas requiere planificación en lenguaje y los productos del trabajo son, en gran medida, lenguaje o contenido estructurado; por eso el LM es hoy el núcleo rector más natural.

\begin{warningbox}{Aquí hay una divergencia valiosa}
Xie Saining enfatiza que el lenguaje no puede sustituir al mundo físico; Zhang Yueguang enfatiza que el lenguaje es el núcleo de inteligencia del Agent. Ambas posturas no se contraponen necesariamente: la inteligencia del mundo físico necesita ir más allá del lenguaje, mientras que los Agent de productividad dependen, en esta etapa, en gran medida del lenguaje para organizar las tareas.
\end{warningbox}

\subsection{Resumen de la sección}

Los productos de Agent no compiten solo en benchmark de tareas largas. Los ciclos cortos, la alta frecuencia, la baja latencia, la retroalimentación del usuario y los resultados editables pueden formar antes una experiencia de producto real. El LM sigue siendo, por ahora, el cerebro central del Agent.
